% Direct-authority diplomatic French for printed p.127.
\SourcePage{127}

\LayerTarget{tohoku:section:1.4}
\noindent\textbf{1.4. Catégories abéliennes.}
On appelle \textit{catégorie abélienne} une catégorie additive $\mathbf{C}$
satisfaisant aux deux axiomes supplémentaires suivants (qui sont autoduals) :

\LayerTarget{tohoku:section:1.4:AB1}
\noindent AB 1) \textit{Tout morphisme admet un noyau et un conoyau}
(cf. 1.3.)

\LayerTarget{tohoku:section:1.4:AB2}
\noindent AB 2) \textit{Soit $u$ un morphisme dans $\mathbf{C}$. Alors le
morphisme canonique
$\bar u:\mathrm{Coim}\,u\to\mathrm{Im}\,u$ \textup{(cf. 1.3)} est un
isomorphisme.}

\LayerTarget{tohoku:section:1.4:p1}
Il en résulte en particulier qu'\textit{une bijection est un isomorphisme}.
Notons qu'il existe de nombreuses catégories additives satisfaisant à AB 1) et
où les morphismes $\bar u:\mathrm{Coim}\,u\to\mathrm{Im}\,u$ sont
toujours bijectifs, sans être nécessairement des isomorphismes. Il en est ainsi
par exemple de la catégorie additive des modules topologiques séparés sur un
anneau topologique donné, en y prenant comme morphismes les homomorphismes
continus, ou de la catégorie des groupes abélien \textit{filtrés}. Autre
exemple moins évident : la catégorie additive des espaces fibrés holomorphes à
fibre vectorielle sur une variété holomorphe de dimension complexe 1. Ce sont
donc là des catégories additives non abéliennes.

\LayerTarget{tohoku:section:1.4:p2}
Si $\mathbf{C}$ est une catégorie abélienne, alors tout le formalisme habituel des
diagrammes d'homomorphismes entre groupes abéliens peut se développer à
nouveau en y remplaçant les homomorphismes par des morphismes dans $\mathbf{C}$, pour
autant qu'on n'a en vue que des propriétés ``de caractère fini'', i. e. ne
faisant pas intervenir des sommes directes ou produits directs infinis (pour
lesquels des précautions spéciales sont nécessaires, voir
N\textsuperscript{o}5). Nous nous contentons ici d'indiquer quelques faits
particulièrement importants, renvoyant pour d'autres détails à [3].

\LayerTarget{tohoku:section:1.4:p3}
Dans la suite, nous nous plaçons dans une catégorie abélienne fixée $\mathbf{C}$. Soit
$A\in\mathbf{C}$, et à tout sous-truc de $A$ faisons correspondre son conoyau (qui est
donc un quotient de $A$), et à tout truc quotient de $A$ faisons correspondre
son noyau (qui est donc un sous-truc de $A$). On obtient ainsi \textit{une
correspondance biunivoque entre la classe des sous-trucs de $A$ et la classe
des trucs quotients de $A$}, cette correspondance étant d'ailleurs un
\textit{antiisomorphisme} pour les relations d'ordre naturelles. D'ailleurs,
les sous-trucs de $A$ forment une classe \textit{réticulée} (donc aussi les
trucs quotients) : Si $P$ et $Q$ sont deux sous-trucs de $A$, leur sup est
l'image de la somme directe $P+Q$ par le morphisme dont les composantes sont
les injections canoniques de $P$ et $Q$ dans $A$, et leur inf est le noyau du
morphisme de $A$ dans le produit $(A/P)\times(A/Q)$ dont les composantes sont les surjections canoniques de $A$ sur
$A/P$ et $A/Q$. (Conformément à l'usage, on désigne par $A/P$ le quotient de
$A$ qui correspond au sous-truc $P$ ; une notation duale comme
$A\backslash R$ pour le sous-truc de $A$ qui correspond à un truc quotient $R$
semble naturelle). On a des interprétations duales pour le inf ou le sup de
deux trucs quotient de $A$.

\LayerTarget{tohoku:section:1.4:p4}
Soit $u:A\to B$ un morphisme. Si $A'$ est un sous-truc de $A$, on
définit l'image de $A'$ par $u$, notée $u(A')$, comme $\mathrm{Im}\,ui$, où
$i$ est l'injection canonique $A'\to A$. Dualement, on définit
l'image inverse $u^{-1}(B')$ d'un quotient $B'$ de $B$, c'est un quotient de
$A$. Si $B'$ est maintenant un sous-truc de $B$, on définit l'image inverse de
$B'$ par $u$, notée $u^{-1}(B')$, comme le noyau de $ju$, où $j$ est

\TohokuContinuation{128}
